\documentclass{article}
\title{Chunking Strategies in AutoRAG}
\author{AutoRAG Documentation}
\date{2026}

\begin{document}
\maketitle

\section{Introduction}
Document chunking is a critical step in the RAG pipeline. The choice of chunking strategy directly impacts retrieval quality and the overall performance of the RAG system. AutoRAG supports two chunking methods: hybrid and recursive.

\section{Hybrid Chunking}
The hybrid chunking method uses Docling's HybridChunker, which is structure-aware. It preserves document hierarchy by respecting heading boundaries, table cells, and figure captions. Chunks are enriched with contextual information from parent headings, improving retrieval relevance.

Hybrid chunking does not support chunk overlap. Each chunk boundary aligns with a structural element in the document. The maximum chunk size is configurable between 128 and 2048 tokens.

\section{Recursive Chunking}
The recursive chunking method uses LangChain's RecursiveCharacterTextSplitter. It first exports the DoclingDocument to markdown format, then splits the text recursively using a hierarchy of separators (paragraph breaks, newlines, spaces). This method supports configurable chunk overlap from 0 to 256 tokens.

\section{Token Estimation}
Both chunkers use an approximate tokenizer that estimates one token per four characters. This avoids the overhead of loading a full tokenizer model while providing reasonable chunk size control.

\end{document}